% Slides for the scaled dot-product attention study.
% Build with tools/build_slides.sh studies/2026-08_scaled-dot-product-attention.
\documentclass[aspectratio=169,11pt]{beamer}

\usetheme{Madrid}
\usecolortheme{default}
\setbeamertemplate{navigation symbols}{}
\usepackage[T1]{fontenc}
\usepackage{amsmath,amssymb}
\usepackage[numbers]{natbib}

\title{Why Scaled Dot-Product Attention Divides by $\sqrt{d_k}$}
\author{self-study-with-ai}
\date{\today}

\begin{document}

\begin{frame}
  \titlepage
\end{frame}

\begin{frame}{Question and scope}
  \begin{itemize}
    \item Why does the Transformer scale attention logits by $1/\sqrt{d_k}$?
    \item Scope: what the primary source, Vaswani et~al.~\citep{vaswani2017attention}, states.
    \item Out of scope: attention variants, later empirical work.
  \end{itemize}
\end{frame}

\begin{frame}{The mechanism}
  \begin{equation*}
    \mathrm{Attention}(Q, K, V)
      = \mathrm{softmax}\!\left(\frac{QK^\top}{\sqrt{d_k}}\right) V
  \end{equation*}
  \begin{itemize}
    \item Stated motivation~\citep{vaswani2017attention}: at large $d_k$, dot
      products grow in magnitude and push softmax into regions of extremely
      small gradients; dividing by $\sqrt{d_k}$ counteracts the effect.
    \item Standard reading (an interpretation, not source text): a logit sums
      $d_k$ roughly unit-scale products, so its standard deviation grows as
      $\sqrt{d_k}$; the division normalizes it.
  \end{itemize}
\end{frame}

\begin{frame}{Why bother: the cost angle}
  \begin{itemize}
    \item Dot-product attention reduces to optimized matrix multiplication, so
      it is faster and more space-efficient than additive
      attention~\citep{vaswani2017attention}.
    \item The scaling factor keeps the cheap form competitive at large
      $d_k$~\citep{vaswani2017attention}.
  \end{itemize}
\end{frame}

\begin{frame}{Evidence and limitations}
  \begin{itemize}
    \item Evidence: single primary source, peer-reviewed, NeurIPS
      2017~\citep{vaswani2017attention}.
    \item The motivation is asserted analytically, not measured: no ablation of
      scaled versus unscaled attention appears in the source.
    \item Multi-head configuration: $h = 8$, $d_k = 64$, $d_{\mathrm{model}} =
      512$~\citep{vaswani2017attention}.
  \end{itemize}
\end{frame}

\begin{frame}{References}
  \footnotesize
  \bibliographystyle{plainnat}
  \bibliography{refs}
\end{frame}

\end{document}
